\section{Theory}	\label{a:stil_en}

In this chapter, the theoretical foundations of this work are developed. We begin with Gaussian processes as probabilistic function approximators, 
present Bayesian optimization as a framework for sample efficient black-box optimization and then eventually look over the Cheetah simulator and the 
ARES accelerator where the work is used and how physics informed knowledge may be integrated through the prior mean function.


\subsection{Gaussian Process}

Gaussian process (GP) offers a robust framework for regression which allows uncertainty quantification along with predictions. GPs are well suited for 
modeling unknown objectives in optimization problems \cite{roussel2024bo} because they define probability distributions directly over functions, unlike parametric models 

+that learn a predefined set of weights. A Gaussian process is fully specified by its mean function $m(x)$ and covariance function (kernel) $k(x,x')$ that encodes 
assumptions about the function’s smoothness:
\begin{equation}
f(x) \sim \mathcal{GP}(m(x), k(x,x'))
\end{equation}

The kernel carries assumptions regarding smoothness and correlation between points, the mean function encodes previous expectations regarding the behavior of the function. 
The kernel selection determines what kind of functions the GP considers probable. Because of its ability in modeling smooth, physically realistic functions without assuming 
infinite differentiability,  Mat\'ern-5/2 kernel is widely used in Bayesian Optimization \cite{roussel2024bo}. The GP uses Bayes’s rule (formula for baye’s rule) to update 
from prior to posterior when we observe data. The posterior mean remains Gaussian, with mean and variance that reflects both our prior assumption and observed evidence. 
One of the important characteristics is that posterior uncertainty is low in areas with dense observations and high in areas with sparse data, this natural uncertainty 
quantification is what makes GPs useful for guiding optimization.

\subsection{Bayesian Optimization}
Bayesian Optimization (BO) is a sequential method for determining the optimal value of expensive black-box functions \cite{roussel2024bo}. It is especially effective when gradient information 
is not available and each function evaluation is expensive. Because of these characteristics, BO is ideally suited for accelerator tuning, where it has emerged as a cutting-edge 
method \cite{jalas2021bolaserplasma} \cite{Duris2020bofreeelectron}. The goal is to solve:

\begin{equation}
    \label{eq_optFormulation}
    x^* = \arg\min_{x \in \mathcal{X}} f(x)
\end{equation}

Here, $x$ represents the control variables (like magnet settings), $\mathcal{X}$ is the search space and $f(x)$ is the quantity to be minimised (beam error in our case). BO aims to find $x^*$ using as 
few function evaluation as possible. The key idea is to use a probabilistic surrogate model (usually a GP) to guide the search rather than evaluating the expensive function everywhere. \

% As shown in Figure~\ref{fig:bayesian_optimization_loop}, the optimization process is iterative, We start with a modest initial set of observations, fit a GP model to the data and then choose the next point to evaluate using an acquisition 
% NOTE: The reference fig:bayesian_optimization_loop is missing or not defined. Uncomment and fix the label if the figure is present in the main document.
function. Exploration (sampling when uncertainty is large) and exploitation (sampling where the predicted value is good) are balanced by the acquisition function. After evaluating the true objective at the 
selected point, we update the dataset and repeat until the evaluation budget is exhausted.

\begin{figure}[ht]
\centering
\begin{tikzpicture}[
    node distance=1.0cm,
    block/.style={
        rectangle, 
        draw=black!60, 
        fill=black!5, 
        thick,
        minimum width=6cm, 
        minimum height=1.2cm, 
        text centered, 
        rounded corners=4pt,
        drop shadow={shadow xshift=1pt, shadow yshift=-1pt, black!30!gray}
    },
    arrow/.style={
        ->, 
        >=stealth, 
        thick,
        black!80!black
    },
    repeat/.style={
        font=\bfseries\small,
        color=black!70!black
    }
]

% Nodes (vertical flow)
% Replace all the problematic node lines with this:

\node[block] (init) {\parbox{5.5cm}{\centering\textbf{Initialisation} \\ (Random samples)}};
\node[block, below=of init] (gp) {\parbox{5.5cm}{\centering\textbf{Fit GP Model} \\ $f \sim GP(m, k)$}};
\node[block, below=of gp] (acq) {\parbox{5.5cm}{\centering\textbf{Optimise Acquisition} \\ $x_{t+1} = \arg\max \alpha$}};
\node[block, below=of acq] (eval) {\parbox{5.5cm}{\centering\textbf{Evaluate Objective} \\ $y_{t+1} = f(x_{t+1})$}};
\node[block, below=of eval] (update) {\parbox{5.5cm}{\centering\textbf{Update Dataset} \\ $D_{t+1} = D_t \cup \{\text{new}\}$}};

% Main flow arrows
\draw[arrow] (init) -- (gp);
\draw[arrow] (gp) -- (acq);
\draw[arrow] (acq) -- (eval);
\draw[arrow] (eval) -- (update);

% Loop back arrow with "Repeat" label
\draw[arrow, black!70!black, thick] 
    (update.south) -- ++(0,-0.8) 
    -| ++(-3.8,0) 
    |- node[repeat, pos=0.2, left] {Repeat} ([xshift=-0.3cm]gp.west);

\end{tikzpicture}
\caption{The Bayesian optimisation loop.}
\label{fig:bayesian_optimization_loop}
\end{figure}


In this project work, the predicted mean and uncertainty are directly combined using the Upper Confidence Bound (UCB) acquisition function:
\begin{equation}
    \alpha_{UCB}(x) = \mu(x) + \beta \sigma(x)
\end{equation}

where $\mu(x)$ and $\sigma(x)$ are the GP posterior mean and standard deviation respectively. The parameter $\beta$ controls the exploration-exploitation trade-off. 
For minimization, this formulation prefers points that either have low predicted values (exploitation) or high uncertainty (exploration). We use $\beta$ = 2.0 throughout this project work, following the common practice \cite{roussel2024bo} .

\subsection{Physics Informed Prior Mean Function}

Before any data is observed, our belief about the objective is represented by the prior mean function $\mu(x)$. The standard practice is to use zero or constant mean, which assumes no prior knowledge. 
However, physics models that approximate the system are often used in scientific applications, incorporating this knowledge can significantly improve the performance \cite{boltz2025priormeanmodule} . \

To understand why, consider the GP posterior mean prediction at a new point. The GP learns to predict the difference between the prior mean and the actual function. These residuals are small and can be 
learned from fewer observations when the prior mean is accurate. Whereas when the prior is simply zero, the GP has to learn the entire function from scratch. This motivates using a physics informed similar as the prior mean:

\begin{equation}
    m(x) = f_{physics}(x; \theta)
\end{equation}

Where $\theta$ represents the model parameters like element alignments or beam characteristics. This benefits to improve sample efficiency, have meaningful predictions before observing data and physical consistency in extrapolation.\

A key innovation in this work is making the physics model parameters $\theta$ trainable alongside the standard GP hyperparameters. In standard GP practice, the kernel hyperparameters $\omega$ (length scales, signal variance) and 
noise variance are estimated by maximising the log marginal likelihood of observed data \cite{williams2006gaussian} . As theeta influences the prior mean predictions, it can be included in this optimisation. 


\begin{equation}
    \theta^* = \arg\max_\theta \log p(y | X, \theta, \omega)
\end{equation}

For this optimization to work, the physics similar must be differentiable so that we can compute how changes in the model parameters $\theta$ affect the marginal likelihood. Cheetah satisfies this requirement through its PyTorch implementation. 
As the optimization progresses and more data is collected, the GP fitting process automatically adjusts the physics model parameters to better match the observed beam behavior. This is the type of online system identification where the model 
learns hidden system parameters as a byproduct of optimization itself, eliminating the need for calibration experiments. In this project, we use this approach to learn quadrupole misalignments  which are small transverse offsets that effect 
the beam steering but cannot be directly measured during accelerator operation.



\subsection{Cheetah Based Differential Simulation Library}

Using physics knowledge based simulator as a prior mean function requires two key properties: the simulator must be fast enough to call within each GP evaluation, and the simulator must be differentiable in order to facilitate gradient based 
hyperparameter optimization. A beam dynamics code called Cheetah \cite{kaiser2024bridging} was created to satisfy these needs.  Cheetah uses transfer matrix concept to depict beam transport. For linear optics, each accelerator element transforms 
the beam’s phase space coordinates through a 7x7 transfer matrix. The standard 6x6 transfer matrix Ro is embedded within this structure, with the additional dimension enabling representation of effects such as magnet misalignments. 
When compared to conventional simulation codes, Cheetah achieves speedups of several orders of magnitude through optimized tensor computation and automatic transfer map reduction. Moreover, all operations in Cheetah is implemented entirely 
on PyTorch, which means that all operations in Cheetah support automatic differentiation. Gradients of beam properties with respect to any element parameter including magnet position (misalignments) and strengths can be computed efficiently. 
This ability of differentiation is important for our approach, when GP fits its hyperparameter by maximising marginal likelihood, gradients flow through the Cheetah simulation to update the trainable misalignment parameter. Cheetah provides 
two beam representations. 

\begin{itemize}
    \item \textbf{Parameter Beam:} assumes a Gaussian beam and describes it using a seven dimensional mean vector and covariance matrix, enabling really fast simulation. 
    \item \textbf{Particle Beam:} tracks individual macro-particles for higher fidelity results but then there is a tradeoff of higher computational cost. 
\end{itemize}


Since the ARES experimental area runs in the linear regime, where moment-based tracking is sufficiently accurate, this project work uses Parameter Beam.


\subsection{The ARES accelerator}

The Accelerator Research Experiment at SINBAD (ARES) facility at DESY serves as a testbed for advanced accelerator concepts and to test machine learning 
techniques on them. It is a normal-conducting S-band linac which produces electron beams up to 155 MeV and is designed for flexibility and reproducibility 
in research applications. \

The project work focuses on the ARES experimental area, as demonstrated in figure 2.2. The beamline consists of three quadrupole magnets (Q1, Q2, Q3) and 
two corrector/steering magnets (CV for vertical steering and CH for horizontal steering), and a diagnostic screen where beam properties are measured. 

The beam tuning task involves optimizing these five magnets:  
\begin{equation}
    \mathbf{u} = (k_{Q1}, k_{Q2}, k_{Q3}, \theta_{CV}, \theta_{CH})
\end{equation}
where $\mathbf{u} \in \mathbb{R}^5$ is the control vector. \

In order to produce a focused and centred beam at screen, The measured observables are beam’s 
horizontal and vertical positions ($\mu_x$ and $\mu_y$) and sizes ($\sigma_x$ and $\sigma_y$). The objective 
function used throughout this work is the mean absolute error which combines both centering and 
focusing into a single metric to minimise
\begin{equation}
    MAE = \frac{1}{4}(|\mu_x| + \sigma_x + |\mu_y| + \sigma_y)
\end{equation}

Partial observability one of the major challenge in this problem, the quadrupole magnets are never exactly aligned; 
the precise beam distribution entering the experimental area is not known. Even with the corrector magnets set to zero, 
a misaligned quadrupole with transverse offset ($\Delta_x$, $\Delta_y$) partially functions as steering element eventually 
diverting the beam. These misalignments are not directly measurable during operation. \ 

This Simulation-to-Reality gap shows that even an accurate physics model will not perfectly predict beam behavior unless it 
counts for the actual misalignments. By treating the six misalignments parameters (horizontal and vertical offsets for each quadrupole) 
as trainable quantities within the physics-informed prior, our approach addresses this issue and allows the model to adjust to actual 
accelerator settings using the optimization data. 





% \usepackage{tikz}
% \usetikzlibrary{shapes, arrows, positioning, shadows}


\begin{figure}[ht]
\centering
\begin{tikzpicture}[
    node distance=2.5cm,
    magnet/.style={
        rectangle,
        draw=blue!50,
        fill=blue!10,
        thick,
        minimum width=1.2cm,
        minimum height=0.8cm,
        text centered
    },
    screen/.style={
        rectangle,
        draw=green!50,
        fill=green!10,
        thick,
        minimum width=2cm,
        minimum height=1.5cm,
        text centered,
        rounded corners=4pt
    },
    beam/.style={
        -latex,
        thick,
        red!70!black,
        line width=1.5pt
    },
    label/.style={
        font=\small\bfseries
    },
    param/.style={
        font=\scriptsize,
        color=blue!70!black
    },
    elemname/.style={
        font=\tiny,
        color=black!70
    }
]

% Beam line
\draw[beam] (0,0) -- (10,0);
\node at (0,-0.3) [label] {beam};

% Quadrupole Q1
\node[magnet] (Q1) at (1.5,0) {Q1};
\node[above=0.2cm of Q1, param] {$k_1$};
\node[below=0.4cm of Q1, elemname] {AREAMQZM1};

% Quadrupole Q2
\node[magnet] (Q2) at (3.5,0) {Q2};
\node[above=0.2cm of Q2, param] {$k_1$};
\node[below=0.4cm of Q2, elemname] {AREAMQZM2};

% Corrector CV
\node[magnet] (CV) at (5.5,0) {CV};
\node[above=0.2cm of CV, param] {$\theta_v$};
\node[below=0.4cm of CV, elemname] {AREAMCVM1};

% Quadrupole Q3
\node[magnet] (Q3) at (7.5,0) {Q3};
\node[above=0.2cm of Q3, param] {$k_1$};
\node[below=0.4cm of Q3, elemname] {AREAMQZM3};

% Corrector CH
\node[magnet] (CH) at (9.5,0) {CH};
\node[above=0.2cm of CH, param] {$\theta_h$};
\node[below=0.4cm of CH, elemname] {AREAMCHM1};

% Screen
\node[screen] (SCREEN) at (12,0) {SCREEN};
\node[below=0.4cm of SCREEN, elemname] {AREABSCR1};

% Arrows between elements
\draw[thick, black!50, ->] (Q1) -- (Q2);
\draw[thick, black!50, ->] (Q2) -- (CV);
\draw[thick, black!50, ->] (CV) -- (Q3);
\draw[thick, black!50, ->] (Q3) -- (CH);
\draw[thick, black!50, ->] (CH) -- (SCREEN);

\end{tikzpicture}
\caption{Schematic of the ARES experimental area lattice showing the magnet sequence. The beam travels from left to right through quadrupoles (Q1, Q2, Q3) and correctors (CV, CH) before reaching the screen. The adjustable parameters are the quadrupole strengths $k_1$ and corrector angles $\theta_v$, $\theta_h$.}
\label{fig:ares_lattice}
\end{figure}




